\documentclass[12pt]{article}
\usepackage[margin=1in,a4paper]{geometry}
\usepackage{amsmath,amssymb,fullpage,xcolor,mhchem,mdframed}
\usepackage{hyperref}
\DeclareMathOperator\Tr{Tr}
\begin{document}
\title{Chapter 43: Majorana Fields}
\date{\today}
\maketitle

\begin{mdframed}[backgroundcolor=blue!5]
\textbf{Ch.43 covers a special spinor: the Majorana field,}
where the particle is its own antiparticle.  Appears in SUSY (gauginos,
higgsinos) and neutralino dark matter.
\end{mdframed}

\section{Majorana Condition}
\begin{mdframed}
$\displaystyle\psi(x) = \psi^c(x) \equiv C\,\bar\psi^T(x)$ ,
\quad $C=i\gamma^2\gamma^0$, $C^2=-\mathbf{1}$, $C\gamma^\mu C^{-1}=-(\gamma^\mu)^T$ .
\end{mdframed}
This halves the DOF: 2 (like Weyl) not 4 (like Dirac).

\section{Comparison}
\begin{center}
\begin{tabular}{|l|c|c|c|}
\hline
Property & Weyl & Majorana & Dirac \\
\hline
DOF & 2 & 2 & 4 \\
EM charge & 0 & 0 & $\pm e$ \\
Lorentz rep & $(\tfrac12,0)$ & Singlet & $(\tfrac12,0)\oplus(0,\tfrac12)$ \\
SUSY partner & $\tilde\psi$ & Gaugino & $\tilde f_L+\tilde f_R$ \\
\hline
\end{tabular}
\end{center}

\section{Majorana Propagator}
\begin{mdframed}
$S_F(k) = i/(k\!\!\!/+m)$.  But Wick contraction gives $\tfrac12 S_F$ (extra $\tfrac12$ factor).
\end{mdframed}
This $\tfrac12$ is the statistical origin of Majorana statistics.

\section{Feynman Rules}
\begin{center}
\begin{tabular}{|l|c|}
\hline
Object & Rule \\
\hline
Majorana propagator & $i/(k\!\!\!/+m)$ \\
Majorana Yukawa vertex & $-iy\,\gamma^\mu$ \\
Closed Majorana loop & $\frac12\Tr[S_F]$ \\
4-Majorana contact & $-iy^2$ \\
\hline
\end{tabular}
\end{center}

\section{Majorana Mass}
$\mathcal{L}\supset -\tfrac12 m_M\,\psi^TC\psi + \text{h.c.}$
The $\tfrac12$ reflects: only half the DOF (particle=antiparticle).

MSSM examples: $\tfrac12 M_1\tilde\lambda\tilde\lambda$ (gaugino Majorana mass).

\section{Why Majorana Matters for MHV}
In SUSY extensions of QCD, gluinos are Majorana fermions.
They enable double-pole signatures at the LHC.
The BCFW recursion generalizes to SUSY theories using massless
helicity states: the Majorana property is not an obstacle but a feature.

\vfill\bibliographystyle{plain}\end{document}